% TEMPLATE-MILC-v3.tex - plantilla maestra para todas las guias MILC v3
% Ver scripts/generadores/build-guia-milc.py para la lista completa de
% claves esperadas. El script valida que no queden placeholders sin
% reemplazar antes de escribir el archivo final.

\documentclass[11pt,letterpaper]{article}
\usepackage[margin=1.75cm]{geometry}
\usepackage[table]{xcolor}
\usepackage{fontspec}
\usepackage[spanish]{babel}
\usepackage{tikz}
\usepackage[most]{tcolorbox}
\usepackage{tabularx}
\usepackage{array}
\usepackage{fancyhdr}
\usepackage{titlesec}
\usepackage{ragged2e}
\usepackage{enumitem}
\usepackage{microtype}

\setmainfont{Helvetica Neue}
\setsansfont{Helvetica Neue}
\setlength{\parindent}{0pt}
\setlength{\parskip}{6pt}
\hyphenpenalty=200
\exhyphenpenalty=200
\pretolerance=400
\tolerance=2000
\setlength{\emergencystretch}{1em}
\renewcommand{\arraystretch}{1.25}
\setlist{leftmargin=5mm,itemsep=1mm,topsep=1mm}
\newcolumntype{Y}{>{\RaggedRight\arraybackslash}X}

\definecolor{milcMagenta}{HTML}{E6007E}
\definecolor{milcVino}{HTML}{5A0038}
\definecolor{milcTurquesa}{HTML}{0093A5}
\definecolor{milcVerde}{HTML}{58A923}
\definecolor{milcMostaza}{HTML}{E5B400}
\definecolor{milcOcre}{HTML}{B86600}
\definecolor{milcNegro}{HTML}{241816}
\definecolor{milcGris}{HTML}{F7F7F4}
\definecolor{milcLinea}{HTML}{E8E2DD}
\definecolor{milcGradePrimary}{HTML}{0066FF}
\definecolor{milcGradeDark}{HTML}{003D99}
\definecolor{milcGradeSoft}{HTML}{D6E8FF}
\definecolor{milcGradeAccent}{HTML}{CFE7FF}
\definecolor{milcGradeText}{HTML}{FFFFFF}
\color{milcNegro}

\pagestyle{fancy}
\fancyhf{}
\lhead{\small Tecnología e Informática · Grado 11}
\rhead{\small Guía 24 · MILC v3}
\cfoot{\small\thepage}
\renewcommand{\headrulewidth}{0pt}

\newtcolorbox{titlebox}[1]{
  enhanced,colback=#1,colframe=#1,arc=7mm,boxrule=0pt,
  left=7mm,right=7mm,top=3mm,bottom=3mm,
  before skip=8pt,after skip=8pt,
  fontupper=\sffamily\bfseries\Large\color{white}
}

\newtcolorbox{softbox}[3]{
  enhanced,colback=#2,colframe=#1,borderline west={4pt}{0pt}{#1},
  arc=2mm,boxrule=.4pt,left=4mm,right=4mm,top=3mm,bottom=3mm,
  before skip=6pt,after skip=8pt,title={#3},coltitle=white,
  colbacktitle=#1,fonttitle=\sffamily\bfseries,fontupper=\RaggedRight,
  attach boxed title to top left={xshift=3mm,yshift=-2mm},
  boxed title style={arc=2mm, boxrule=0pt}
}

\newtcolorbox{drawbox}[2]{
  enhanced,colback=white,colframe=#1,arc=4mm,boxrule=1pt,
  left=5mm,right=5mm,top=4mm,bottom=4mm,before skip=8pt,after skip=8pt,
  title={#2},coltitle=white,colbacktitle=#1,fonttitle=\sffamily\bfseries,
  fontupper=\RaggedRight,
  attach boxed title to top left={xshift=4mm,yshift=-2mm},
  boxed title style={arc=3mm, boxrule=0pt}
}

\newtcolorbox{infoband}[1]{
  enhanced,colback=#1!8,colframe=#1!50,arc=2mm,boxrule=.5pt,
  left=4mm,right=4mm,top=2mm,bottom=2mm,before skip=4pt,after skip=4pt,
  fontupper=\small\RaggedRight
}

% Puente narrativo entre secciones (contrato editorial Conéctate).
% Da continuidad: "Acabas de X. Ahora vas a Y porque Z."
% Visualmente: banda izquierda en color del grado, texto en itálica pequeña.
\newtcolorbox{puentebox}{
  enhanced,colback=milcGradeSoft,colframe=milcGradeDark,
  borderline west={3pt}{0pt}{milcGradeDark},
  arc=0mm,boxrule=0pt,left=5mm,right=4mm,top=2.5mm,bottom=2.5mm,
  before skip=4pt,after skip=8pt,
  fontupper=\itshape\small\color{milcNegro}\RaggedRight
}
\newcommand{\puente}[1]{%
  \begin{puentebox}%
    \textcolor{milcGradeDark}{\textbf{\textrightarrow}}\hspace{2mm}#1%
  \end{puentebox}%
}

\newcommand{\linead}[1]{\noindent\color{milcLinea}\rule{#1}{0.5pt}\color{milcNegro}}
\newcommand{\respuesta}[1]{\par\vspace{2mm}\foreach \n in {1,...,#1}{\noindent\color{milcLinea}\rule{\linewidth}{0.5pt}\par\vspace{5mm}}\color{milcNegro}}
\newcommand{\checkbox}{\textcolor{milcGradeDark}{\fbox{\phantom{x}}}\hspace{5pt}}

\newcommand{\phasecard}[4]{
\begin{tcolorbox}[enhanced,colback=#3,colframe=#2,arc=3mm,boxrule=.5pt,height=3.05cm,valign=center,halign=center,left=1.5mm,right=1.5mm,top=2mm,bottom=2mm]
{\sffamily\bfseries\fontsize{22}{24}\selectfont\textcolor{#2}{#1}}\par
{\sffamily\bfseries\small #4}
\end{tcolorbox}
}

\begin{document}
\thispagestyle{empty}
\begin{tikzpicture}[remember picture,overlay]
  \fill[white] (current page.south west) rectangle (current page.north east);
  \path[fill=milcGradePrimary,rounded corners=16mm]
    ([xshift=.55cm,yshift=-.55cm]current page.north west) rectangle
    ([xshift=-.55cm,yshift=3.65cm]current page.south east);

  \node[anchor=north west,text=milcGradeText] at ([xshift=2.0cm,yshift=-1.85cm]current page.north west)
    {\sffamily\bfseries\fontsize{14}{16}\selectfont GRADO};
  \node[anchor=north west,text=milcGradeText,opacity=.82] at ([xshift=2.0cm,yshift=-2.75cm]current page.north west)
    {\sffamily\bfseries\fontsize{9}{11}\selectfont SERIE GUÍAS · TECNOLOGÍA E INFORMÁTICA};

  \begin{scope}[shift={([xshift=-3.05cm,yshift=-2.55cm]current page.north east)}]
    \fill[white,opacity=.80] (-.62,-.52) rectangle (.62,.52);
    \draw[milcGradeText,line width=1pt] (-.62,-.52) rectangle (.62,.52);
    \fill[milcGradeDark] (-.42,-.38) rectangle (-.22,.06);
    \fill[milcGradeAccent] (-.10,-.38) rectangle (.10,.24);
    \fill[milcGradePrimary!60!black] (.22,-.38) rectangle (.42,.38);
  \end{scope}

  \node[anchor=west,text=milcGradeText] at ([xshift=1.9cm,yshift=-12.85cm]current page.north west)
    {\sffamily\bfseries\fontsize{124}{124}\selectfont 11};
  \draw[milcGradeText,line width=1.7pt]
    ([xshift=4.52cm,yshift=-11.68cm]current page.north west) circle (.13cm);

  \node[anchor=west,text=milcGradeText,text width=8.7cm,align=left] at ([xshift=10.35cm,yshift=-11.6cm]current page.north west)
    {
     {\sffamily\bfseries\fontsize{19}{22}\selectfont MVP digital ---\\una versión imperfecta\\para probar}\\[4mm]
     {\sffamily\bfseries\fontsize{23}{26}\selectfont Guía 24-11-TIC}\\[2mm]
     {\sffamily\fontsize{13}{16}\selectfont Período 3 · Proyecto final emprendedor}\\[5mm]
     {\sffamily\bfseries\fontsize{11}{13}\selectfont Institución Educativa Sor María Juliana}\\[1mm]
     {\sffamily\fontsize{10}{12}\selectfont Autor: Dr. Álvaro Cárdenas Orozco}
    };

  \path[fill=milcGradeSoft,rounded corners=7mm]
    ([xshift=1.35cm,yshift=.85cm]current page.south west) rectangle
    ([xshift=-1.35cm,yshift=4.25cm]current page.south east);
  \draw[milcGradeDark,line width=.65pt,rounded corners=7mm]
    ([xshift=1.35cm,yshift=.85cm]current page.south west) rectangle
    ([xshift=-1.35cm,yshift=4.25cm]current page.south east);
  \node[anchor=center,text=milcNegro,text width=17.0cm,align=center] at ([yshift=2.55cm]current page.south)
    {\sffamily\fontsize{10}{12}\selectfont Metodología MILC · Apertura ancestral · Pensamiento computacional · Triángulo Dussel-Estoicismo-Floridi\\
    \textcolor{milcGradeDark}{\bfseries Producto:} MVP funcional (Producto Mínimo Viable) que combine al menos 3 piezas digitales (sitio simple o landing + formulario de captura + flujo de respuesta) y que pueda ser probado por al menos 1 usuario real esta misma semana, con criterios explícitos de qué se considera \emph{``funcionó''} y qué se considera \emph{``falló''} antes de construirlo};
\end{tikzpicture}
\mbox{}
\newpage

\section*{Apertura: MVP digital --- una versión imperfecta para probar}

\begin{softbox}{milcGradeDark}{milcGradeSoft}{Saber ancestral}
En los talleres alfareros de Cartago la Vieja, en La Tebaida, en el Quindío profundo, ningún maestro saca jamás una serie de 50 piezas sin antes hacer \textbf{la pieza de muestra}. Esa pieza tiene un nombre antiguo: \emph{la primera}. La primera es deliberadamente imperfecta: el barro elegido es del montón principal, pero el horneado se hace a temperatura menor para verificar antes de invertir leña; la forma se hace rápida para validar la mano; el esmalte se aplica solo en una zona para ver cómo se comporta. Si la primera quiebra, no se quema la serie. Si la primera se ve sólida, recién entonces el maestro invierte 3 días en producir las 50. La sabiduría es ancestral y económica: \textbf{nadie quema 50 piezas para descubrir que el barro no servía}. La misma lógica opera con la \textbf{partera} que prueba un té antes de prescribirlo, con el \textbf{ebanista} que entalla una muestra antes del juego de comedor, con el \textbf{sastre} que arma el patrón en papel antes de cortar la tela buena. El oficio profesional siempre ha sabido que \emph{una versión imperfecta para probar} ahorra meses de error.
\end{softbox}

\begin{softbox}{milcTurquesa}{milcGris}{Saber contemporáneo}
El \textbf{MVP, Producto Mínimo Viable} (Eric Ries, \emph{The Lean Startup}, 2011) es la versión digital exacta de la \emph{primera} del alfarero: la versión más pequeña posible del producto que permita aprender algo real sobre los supuestos críticos del Canvas. La regla profesional es famosa y dolorosa: \textbf{si no te avergüenza tu MVP, lo lanzaste demasiado tarde} (Reid Hoffman, fundador de LinkedIn). Un MVP digital tiene 3 propiedades irrenunciables: (1) \textbf{Funciona end-to-end}: un usuario puede recorrer la propuesta de valor de inicio a fin, aunque sea con costuras visibles. (2) \textbf{Verifica un supuesto crítico}: cada MVP está diseñado para probar 1 o 2 de los 3 supuestos críticos identificados en la sesión 3. (3) \textbf{Se construye en días, no en meses}: si tu MVP toma más de 2 semanas en estar listo, no es mínimo. Las herramientas modernas (Google Sites, Tally, Carrd, Notion, WhatsApp Business, Zapier) permiten construir MVP funcionales en una tarde sin escribir código.
\end{softbox}

\begin{softbox}{milcMostaza}{milcGris}{Pregunta-puente}
\textit{¿Qué sabía el alfarero al hacer \emph{la primera} antes de la serie, que el emprendedor novato olvida cuando se obsesiona con \emph{``lanzar perfecto''}? ¿Y por qué Reid Hoffman dice que la vergüenza por tu MVP es señal de buena disciplina, no de mal trabajo?}
\end{softbox}

\begin{softbox}{milcVerde}{milcGris}{Saber y saber hacer}
Al terminar podrás: (1) \textbf{identificar} cuál de tus 3 supuestos críticos vale más la pena probar primero y por qué; (2) \textbf{aplicar} el patrón MVP combinando 3 piezas digitales accesibles (landing + formulario + flujo de respuesta); (3) \textbf{crear} tu MVP funcional, end-to-end, en una sola tarde de trabajo concentrado; (4) \textbf{evaluar} tu MVP con criterios explícitos de éxito y fracaso escritos \emph{antes} de mostrárselo al primer usuario.
\end{softbox}

\small
\begin{tabularx}{\textwidth}{>{\bfseries}p{3.0cm}Y}
\rowcolor{milcGradePrimary!12}Autor & Dr. Álvaro Cárdenas Orozco\\
DBA & Diseño, implementación y sustentación de un proyecto de impacto local (MEN, T\&I)\\
Referentes & Dussel (analéctica · sur global) · Lean Startup · Floridi · Estoicismo\\
Producto final & MVP funcional (Producto Mínimo Viable) que combine al menos 3 piezas digitales (sitio simple o landing + formulario de captura + flujo de respuesta) y que pueda ser probado por al menos 1 usuario real esta misma semana, con criterios explícitos de qué se considera \emph{``funcionó''} y qué se considera \emph{``falló''} antes de construirlo\\
\end{tabularx}
\normalsize

\puente{Antes de empezar, mira el viaje completo. En la sesión 3 cerraste el Canvas con 3 supuestos críticos. Hoy construyes el \textbf{primer artefacto} que pondrá uno de esos supuestos contra el mundo. La sesión 5 medirá lo que pase con 10 usuarios reales; la sesión 6 contará la historia en pitch. Sin MVP funcional hoy, las dos sesiones siguientes no tienen materia prima.}

\begin{titlebox}{milcGradeDark}
Ruta MILC: así vamos a aprender
\end{titlebox}
\begin{tabularx}{\textwidth}{>{\centering\arraybackslash}X>{\centering\arraybackslash}X>{\centering\arraybackslash}X>{\centering\arraybackslash}X}
\phasecard{1}{milcVerde}{milcGris}{Escuta\\[-1mm]\small Observo, escucho y pregunto.} &
\phasecard{2}{milcTurquesa}{milcGris}{Sistematización\\[-1mm]\small Pensamiento computacional.} &
\phasecard{3}{milcMagenta}{milcGris}{Praxis\\[-1mm]\small Construyo y aplico.} &
\phasecard{4}{milcVino}{milcGris}{Evaluación\\[-1mm]\small Reflexión liberadora.}
\end{tabularx}


\puente{Antes de teorizar, vas a hacer una \emph{lectura honesta} de tus 3 supuestos críticos: ¿cuál es el más \textbf{caro de equivocarse}?, ¿cuál es el más \textbf{rápido de verificar}?, ¿cuál es el más \textbf{barato de construir como MVP}? Elegirás uno aplicando la regla de \emph{máximo aprendizaje por hora invertida}.}

\begin{titlebox}{milcVerde}
Fase 1 · Escuta
\end{titlebox}

\begin{softbox}{milcVerde}{milcGris}{Escena para escuchar}
\textbf{Actividad 1 · IDENTIFICA --- ¿Cuál de mis 3 supuestos críticos pruebo primero?} (15 min · individual). Toma la tabla de la sesión 3 con tus 3 supuestos críticos formulados como afirmaciones falsables. Para cada supuesto, evalúa 3 dimensiones en escala 1-5: (a) \textbf{Costo de equivocarse}: si este supuesto fuera falso y yo no lo supiera hasta el final, ¿qué tan grave sería? (b) \textbf{Rapidez de verificación}: ¿en cuántos días puedo probarlo con un MVP simple? (c) \textbf{Barato de construir}: ¿qué tan accesible es montar un MVP que lo pruebe? Calcula \textbf{Prioridad = Costo × Rapidez × Accesibilidad} (multiplicación de los 3) y elige el supuesto con mayor prioridad como el que probarás hoy.
\end{softbox}

\checkbox Evalué los 3 supuestos críticos en las 3 dimensiones.\\
\checkbox Elegí 1 supuesto para probar con MVP, justificado con números.\\
\checkbox El supuesto está formulado como afirmación falsable (no como pregunta).

\begin{infoband}{milcVerde}
\textbf{Lo que va al cuaderno (Actividad 1 · IDENTIFICA).} Título: «Actividad 1 --- Supuesto a probar con MVP». \textbf{Formato:} tabla 3 filas × 4 columnas (Supuesto~|~Costo~|~Rapidez~|~Accesibilidad), con producto y elección final marcada. Al final: 1 línea con el supuesto elegido formulado como afirmación falsable. \textbf{Sabes que terminaste cuando} tienes claro qué prueba tu MVP, no solo qué hace.
\end{infoband}

\puente{Ya tienes el supuesto que vas a probar. Ahora vas a entender la \textbf{anatomía del MVP digital}: las 3 piezas obligatorias (captura, valor, respuesta), los 5 errores que matan al MVP novato (perfeccionismo, alcance gigante, falta de medición, sin canal de captura, sin criterio de éxito), y los \textbf{4 tipos de MVP} reconocibles en el mercado actual.}

\begin{titlebox}{milcTurquesa}
Fase 2 · Sistematización con pensamiento computacional
\end{titlebox}

\begin{softbox}{milcTurquesa}{milcGris}{Cuatro pilares aplicados al tema}
Un MVP novato es un sitio bonito que no prueba nada. Un MVP profesional es un dispositivo de aprendizaje barato y rápido. Para que el tuyo sea del segundo tipo necesitas la \textbf{anatomía del MVP digital}: las 3 piezas obligatorias, los 4 tipos reconocidos y los errores que matan el aprendizaje antes de empezar.
\end{softbox}

\small
\begin{tabularx}{\textwidth}{>{\bfseries\RaggedRight\arraybackslash}p{3.6cm}Y}
\rowcolor{milcTurquesa!90}\color{white}Pilar & \color{white}\bfseries Aplicación al tema\\
1. Descomposición & Todo MVP digital se descompone en 3 piezas mínimas: (1) \textbf{Captura}: la puerta de entrada donde el usuario \emph{aterriza} (landing en Google Sites, Carrd, Notion público, post de Instagram con link, mensaje pinned de WhatsApp Business). (2) \textbf{Valor}: lo que el usuario \emph{recibe o experimenta} (descarga, formulario para inscripción, video explicativo, demo interactiva, conversación con bot). (3) \textbf{Respuesta}: el flujo que se dispara cuando el usuario actúa (correo automático con Tally, mensaje en WhatsApp, agenda con Calendly, registro en Sheets vía Zapier). Sin los 3, no hay MVP: hay una pieza suelta.\\
2. Reconocimiento de patrones & Existen \textbf{4 tipos de MVP} en el mercado contemporáneo: (a) \textbf{Landing de validación}: página con propuesta y formulario \emph{``quiero esto''}; mide intención. (b) \textbf{Mago de Oz}: la automatización es manual detrás del telón (tú respondes a mano lo que parece automático); mide demanda sin construir tecnología. (c) \textbf{Conserje}: atiendes 1-3 usuarios uno a uno antes de escalar; mide proceso completo a baja escala. (d) \textbf{Pieza única funcional}: construyes solo la función más importante (no el sistema completo); mide si esa función basta para resolver el dolor. Tu MVP elegirá uno de los 4, no los 4.\\
3. Abstracción & Lo esencial cabe en una frase: \emph{el MVP no se construye para vender; se construye para aprender}. Si el resultado del MVP no te enseña algo verdadero sobre tu supuesto crítico, está mal hecho aunque se vea bonito. La phronesis del emprendedor consiste en \textbf{soltar el orgullo del acabado a cambio de la información del experimento}.\\
4. Algoritmo conceptual & Algoritmo: (1) elige el tipo de MVP (de los 4) según tu supuesto; (2) elige 3 herramientas no-code para captura, valor, respuesta; (3) construye end-to-end en una sola sesión de 3-4 horas; (4) prueba tu MVP contigo mismo como primer usuario antes de mostrar a otros; (5) escribe el criterio de éxito \emph{antes} de invitar a un usuario real; (6) prepara la pregunta única que harás al usuario después de usar el MVP: \emph{``¿qué te llevó a hacer X / a no hacer X?''}.\\
\end{tabularx}
\normalsize

\begin{drawbox}{milcTurquesa}{Anatomía del MVP digital}
\textbf{Captura:} dónde aterriza el usuario. \\[1mm] \textbf{Valor:} qué recibe o experimenta. \\[1mm] \textbf{Respuesta:} qué se dispara cuando actúa. \\[1mm] \textbf{Tipo:} landing / Mago de Oz / Conserje / Pieza funcional. \\[1mm] \textbf{Criterio de éxito:} cuántos usuarios hacen qué cosa en cuánto tiempo. \\[1mm] \textbf{Pregunta posterior:} qué le preguntas al usuario después.
\end{drawbox}

\begin{softbox}{milcMostaza}{milcGris}{Errores comunes y su corrección}
(1) \textbf{Perfeccionismo}: gastar 3 semanas puliendo el MVP en vez de lanzarlo en 3 días. (2) \textbf{Alcance gigante}: construir 5 funcionalidades cuando 1 bastaba para probar el supuesto. (3) \textbf{Sin medición}: no definir qué se considera éxito \emph{antes} del lanzamiento. (4) \textbf{Sin canal de captura}: tener MVP listo pero no saber cómo llevar usuarios. (5) \textbf{Sin criterio de fracaso}: no tener escrito \emph{``si esto ocurre, mi hipótesis está muerta''}. (6) \textbf{Construir para impresionar al docente}: el MVP es para el usuario, no para la rúbrica.
\end{softbox}

\begin{infoband}{milcTurquesa}
\textbf{Lo que va al cuaderno (Actividad 2 · APLICA).} Título: «Actividad 2 --- Anatomía de mi MVP». \textbf{Formato:} ficha con los 6 campos de la anatomía aplicados a tu proyecto (captura, valor, respuesta, tipo, criterio de éxito, pregunta posterior). \textbf{Sabes que terminaste cuando} podrías delegar la construcción de tu MVP a otro compañero solo con esa ficha.
\end{infoband}

\puente{Tienes el método. Toca construir. Vas a ensamblar tu MVP en una sola tarde de trabajo combinando 3 herramientas no-code accesibles desde el colegio. Antes de mostrar el MVP a tu primer usuario, escribirás el criterio de éxito explícito: \emph{``este MVP funciona si X usuarios hacen Y dentro de Z días''}.}

\begin{titlebox}{milcMagenta}
Fase 3 · Praxis con pensamiento computacional
\end{titlebox}

\begin{softbox}{milcMagenta}{milcGris}{Construyendo el producto con los cuatro pilares}
\textbf{Actividad 3 · CREA + EVALÚA --- MVP funcional listo para probar} (40 min en sesión, 2-4 horas de construcción durante la semana · individual). Hora de construir. Ensamblas tu MVP combinando 3 herramientas no-code, lo pruebas contigo mismo, escribes el criterio de éxito por adelantado, y agendas al primer usuario real para esta semana.
\end{softbox}

\small
\begin{tabularx}{\textwidth}{>{\bfseries\RaggedRight\arraybackslash}p{3.6cm}Y}
\rowcolor{milcMagenta!90}\color{white}Pilar & \color{white}\bfseries Aplicación al producto\\
Descomposición & Tu MVP se construye en 5 pasos: (1) \textbf{Elige stack no-code}: para captura (Google Sites, Carrd, Notion, Instagram post con link), para valor (Tally, Typeform, video Loom, página con texto, demo con Canva), para respuesta (Zapier, ManyChat, correo automático, WhatsApp Business). (2) \textbf{Conecta las piezas} con 2-3 enlaces simples: captura, valor, respuesta. (3) \textbf{Recorre el MVP tú mismo} simulando ser el usuario; anota cada fricción. (4) \textbf{Arregla las 3 fricciones más graves} (ignora las demás por ahora). (5) \textbf{Genera el link único} que enviarás al primer usuario.\\
Patrones & Sigue el patrón \textbf{``cinta adhesiva sobre código''}: tu MVP es 80\% pegado con cinta adhesiva digital (Zapier conectando 2 apps, copy-paste manual, automatización con un botón humano detrás) y 20\% automatizado. Está bien. Eric Ries llamó esto \emph{``concierge MVP''}: si la operación detrás del telón es manual mientras el frontal se ve automatizado, sigue siendo un MVP válido. La elegancia técnica viene después de validar.\\
Abstracción & El criterio de éxito se escribe \emph{antes} de lanzar y tiene formato fijo: \emph{``Este MVP funcionó si [N] usuarios hacen [acción específica] dentro de [tiempo]''}. Ejemplos reales: \emph{``Si 5 tenderos se inscriben en el formulario en 1 semana''}, \emph{``Si 3 padres de familia agendan demo de Calendly en 3 días''}, \emph{``Si 8 de 10 usuarios completan el formulario sin abandonarlo a mitad''}. Sin número específico y plazo concreto, no hay criterio: hay deseo.\\
Algoritmo de armado & Algoritmo de validación interna del MVP: (1) ¿el usuario entiende la propuesta en menos de 10 segundos al aterrizar? (2) ¿puede hacer la acción principal en menos de 2 minutos? (3) ¿la respuesta llega en menos de 1 hora (automática) o 24 horas (manual)? (4) ¿tienes manera de saber qué hizo el usuario (analytics simple, conteo manual, log en Sheets)? Si las 4 respuestas son sí, tu MVP está listo para el primer usuario real.\\
\end{tabularx}
\normalsize

\begin{drawbox}{milcMagenta}{Checklist antes de invitar al primer usuario}
\checkbox MVP funciona end-to-end (yo lo recorrí completo). \\[1mm] \checkbox 3 piezas conectadas (captura + valor + respuesta). \\[1mm] \checkbox Link único listo para compartir. \\[1mm] \checkbox Criterio de éxito escrito (N usuarios + acción + tiempo). \\[1mm] \checkbox Criterio de fracaso escrito (qué consideraría hipótesis muerta). \\[1mm] \checkbox Manera de saber qué hizo el usuario (sheet, log, analytics). \\[1mm] \checkbox Primer usuario real agendado para esta semana.
\end{drawbox}

\begin{softbox}{milcMagenta}{milcGris}{Plantilla de trabajo}
\textbf{Supuesto a probar.} \textit{Afirmación falsable de la sesión 3.} \\[1mm] \textbf{Tipo de MVP.} \textit{Landing / Mago de Oz / Conserje / Pieza funcional.} \\[1mm] \textbf{Stack.} \textit{Captura + Valor + Respuesta con herramientas concretas.} \\[1mm] \textbf{Link.} \textit{URL única del MVP.} \\[1mm] \textbf{Criterio de éxito.} \textit{N usuarios + acción + tiempo.} \\[1mm] \textbf{Criterio de fracaso.} \textit{Qué resultado mataría la hipótesis.} \\[1mm] \textbf{Primer usuario.} \textit{Nombre + cuándo lo invitas.}
\end{softbox}

\begin{infoband}{milcMagenta}
\textbf{Lo que va al cuaderno (Actividad 3 · CREA + EVALÚA).} Título: «Actividad 3 --- Mi MVP funcional». \textbf{Formato:} ficha con los 7 campos del guion + captura de pantalla impresa o pegada del MVP en cada pieza (captura, valor, respuesta). \textbf{Extensión:} 1-2 páginas. \textbf{Sabes que terminaste cuando} le mandarías el link a un familiar o amigo esta noche sin avergonzarte.
\end{infoband}

\puente{Lo que sigue resume \emph{qué} entregas hoy: MVP funcional + criterio de éxito por escrito + plan para probar con 1 usuario real esta semana. El criterio principal: que un emprendedor experimentado, viendo tu MVP, pueda decir \emph{``entendí qué prueba y cómo se mide''} sin necesidad de explicaciones adicionales.}

\begin{titlebox}{milcMostaza}
Producto de la sesión
\end{titlebox}
\textbf{MVP funcional + criterio de éxito + primer usuario agendado}.
\begin{softbox}{milcMostaza}{milcGris}{Criterios de calidad}
(1) MVP funciona end-to-end con 3 piezas conectadas. (2) Construido en menos de 4 horas con herramientas no-code. (3) Diseñado para probar 1 supuesto crítico específico. (4) Criterio de éxito y de fracaso escritos antes de lanzar. (5) Sistema de medición simple implementado. (6) Primer usuario real agendado para esta semana.
\end{softbox}

\puente{Antes de cerrar, mira el MVP desde las cinco dimensiones humanas. Construir un MVP honesto es respeto por el oficio; construir uno solo para mostrar al docente es desperdicio de la sesión.}

\begin{titlebox}{milcVino}
Fase 4 · Evaluación liberadora — cinco dimensiones
\end{titlebox}

\begin{softbox}{milcVino}{milcGris}{Sentido de la evaluación}
Evaluar no es solo revisar una nota. Es reconocer qué aprendí, qué cambió en mí, qué emociones manejé, cómo me relaciono con la ciudad y la comunidad, y qué le dejaré a quien viene detrás.
\end{softbox}

\textbf{Mi reflexión liberadora en cinco dimensiones.} Completa cada frase iniciada:

\small
\begin{tabularx}{\textwidth}{>{\bfseries\RaggedRight\arraybackslash}p{3.4cm}Y>{\RaggedRight\arraybackslash}p{4.6cm}}
\rowcolor{milcVino}\color{white}Dimensión & \color{white}\bfseries Mi respuesta (frame starter) & \color{white}\bfseries Algo que descubrí\\
1. Personal & Lo que más cambió en mí fue \linead{6.5cm} & \linead{4.2cm}\\
2. Emocional & La emoción más fuerte fue \linead{2.5cm} y la regulé \linead{3cm} & \linead{4.2cm}\\
3. Ciudadana & En democracia esto importa porque \linead{6cm} & \linead{4.2cm}\\
4. Local (Cartago) & En mi barrio sería útil para \linead{6.5cm} & \linead{4.2cm}\\
5. Intergeneracional & A mi abuela/abuelo le diría \linead{6.5cm} & \linead{4.2cm}\\
\end{tabularx}
\normalsize

\begin{infoband}{milcVino}
\textbf{Para la próxima sustentación.} Vuelve a esta tabla en seis meses. Verás que las respuestas cambian — no porque lo aprendido cambie, sino porque tú cambias. La evaluación liberadora es una conversación contigo a través del tiempo.
\end{infoband}


\puente{Y para terminar, tres voces sobre la imperfección útil. Dussel sobre los MVPs que respetan al usuario o lo usan como conejillo de indias, Epicteto sobre la disciplina de soltar la perfección, Floridi sobre la ética del MVP transparente.}

\begin{titlebox}{milcGradeDark}
Triángulo de pensamiento · cierre filosófico
\end{titlebox}

\begin{softbox}{milcVino}{milcGris}{Dussel · lente del nosotros}
\textit{``Todo MVP decide si respeta al usuario como sujeto o lo usa como objeto de experimento.''}

Tu MVP entra en contacto con personas reales. Hay dos formas de pensarlas: como \emph{conejillos de indias} (objetos de prueba para tu aprendizaje) o como \emph{co-investigadores} (sujetos que aceptan ayudarte porque les explicaste el propósito). La filosofía de la liberación pide lo segundo: declara al usuario que es un MVP, que estás aprendiendo, que su retroalimentación es valiosa, que su tiempo no es gratis. La phronesis del oficio digital consiste en \textbf{nunca ocultar al usuario su rol en el experimento}.

\textbf{Cuando invite a mi primer usuario al MVP, ¿le voy a decir con honestidad que es una versión imperfecta para probar, o se lo voy a presentar como producto terminado para impresionar?}
\end{softbox}
\linead{\linewidth}\\[1mm]
\linead{\linewidth}

\begin{softbox}{milcMostaza}{milcGris}{Estoicismo · Epicteto · cuidado interior}
\textit{``Lo que está en tu poder es construir; lo que no está en tu poder es que sea perfecto la primera vez.''}

Es tentador retrasar el lanzamiento del MVP \emph{``hasta que esté mejor''}. La disciplina estoica recomienda lo contrario: \emph{distingue lo que depende de ti (construir y soltar) de lo que no (que sea perfecto)}. El MVP es ejercicio de soltar el control sobre la apariencia para mantener el control sobre el aprendizaje. La perfección es enemiga del oficio cuando ocurre antes que la validación.

\textbf{¿Estoy retrasando mi MVP por mejoras reales o por miedo a que se vea imperfecto? ¿Qué dejaría en el MVP que hoy quiero pulir, si tuviera que lanzarlo en 4 horas?}
\end{softbox}
\linead{\linewidth}\\[1mm]
\linead{\linewidth}

\begin{softbox}{milcTurquesa}{milcGris}{Floridi · lente de la infoesfera}
\textit{``El MVP transparente es la nueva ética del experimento digital con humanos.''}

Antes de internet, los experimentos con personas requerían comités de ética. Hoy cualquier emprendedor lanza MVPs a miles de usuarios sin declarar su naturaleza experimental. La ética digital contemporánea exige un mínimo: el usuario debe saber que está participando en un experimento, qué datos se recogen, qué se hace con ellos. Tu MVP tiene esa obligación a pequeña escala: declara al usuario que es prueba, no producto.

\textbf{¿Mi MVP declara su naturaleza experimental al usuario, o lo disfraza de producto terminado? ¿Estoy entrenándome en una ética que escalaré, o en una que tendré que desaprender?}
\end{softbox}
\linead{\linewidth}\\[1mm]
\linead{\linewidth}

\begin{titlebox}{milcVino}
Compromiso final
\end{titlebox}

\small
\begin{tabularx}{\textwidth}{>{\RaggedRight\arraybackslash}p{3.0cm}YYY}
\rowcolor{milcVino}\color{white}\bfseries Criterio & \color{white}\bfseries Lo logré & \color{white}\bfseries En proceso & \color{white}\bfseries Necesito apoyo\\
Comprensión del tema & Explico ideas centrales con mis palabras. & Reconozco ideas pero falta precisión. & Necesito repasar conceptos base.\\
Pensamiento computacional & Apliqué los 4 pilares al diseñar mi producto. & Apliqué algunos pilares. & Necesito releer la fase 2.\\
Aplicación y producto & Mi producto cumple los criterios y se puede revisar. & Producto funcional con ajustes pendientes. & Aún en construcción.\\
Reflexión 5 dimensiones & Respondí cada dimensión con honestidad. & Respondí algunas con detalle. & Mis respuestas son muy generales.\\
Triángulo de pensamiento & Conecté las 3 voces con mi trabajo real. & Conecté al menos una. & No relaciono las voces con mi proceso.\\
\end{tabularx}
\normalsize

\textbf{Hoy entendí que:}\\[1mm]
\linead{\linewidth}\\[1mm]
\linead{\linewidth}

\textbf{Mi compromiso para la próxima sesión es:}\\[1mm]
\linead{\linewidth}\\[1mm]
\linead{\linewidth}

\begin{softbox}{milcMostaza}{milcGris}{Cita de despedida · Marco Aurelio}
\textit{``El obstáculo en el camino se vuelve el camino. Lo que estorba al camino se vuelve camino.''}\\[1mm]
Cada error de hoy, bien observado, se convierte en pieza del camino que pisarás mañana.
\end{softbox}

\small
\begin{tabularx}{\textwidth}{>{\bfseries}p{3.6cm}Y}
\rowcolor{milcGradeDark!12}Nombre completo: & \linead{10cm}\\
Fecha: & \linead{4cm} \quad Curso/grupo: \linead{4cm}\\
Mi firma: & \linead{9cm}\\
\end{tabularx}
\normalsize

\begin{infoband}{milcGradeDark}
\textbf{Para la próxima sesión.} Llega con tu MVP funcional probado por al menos 1 usuario real y los primeros datos de uso registrados. En la sesión 5 vas a montar el \textbf{plan de medición riguroso} y procesarás los primeros 10 datos. Sin uso real esta semana, la sesión siguiente no tiene materia prima.
\end{infoband}

\end{document}
